\chapter{Generating Functions Meet Arithmetic}

\begin{center}\itshape
Every positive integer was one of his personal friends.\\
\upshape --- J.~E.~Littlewood on Ramanujan, as recalled by G.~H.~Hardy (1921)
\end{center}

In 1918 G.~H.~Hardy and Srinivasa Ramanujan published a formula for $p(n)$, the number of ways of writing $n$ as a sum of
positive integers, order not counting. It was a strange formula, an infinite sum of terms with exponentials of square roots,
and they needed exact values to test it. Major Percy MacMahon, an artillery officer who had become one of the leading
experts on partitions, supplied them. He worked out $p(n)$ for every $n$ up to $200$ by hand. The last entry of his
table was
\[
p(200)=3\,972\,999\,029\,388,
\]
and Hardy and Ramanujan's formula, summed to eight terms, reproduced it with an error of a few thousandths. MacMahon's table
was printed at the end of their paper.

Ramanujan did not stop at checking the formula. He looked at the table the way other people look at faces, and he saw
something nobody had asked about: the numbers $p(4)=5$, $p(9)=30$, $p(14)=135$, $p(19)=490$, $p(24)=1575$, \dots, are all
divisible by $5$. Every fifth entry, starting from the fourth, is a multiple of five, all the way to $200$. There was more:
$p(7n+5)$ is always divisible by $7$, and $p(11n+6)$ by $11$. In a short paper of 1919 he explained that studying
MacMahon's numbers had led him to these ``curious congruence properties'' (we paraphrase his opening), and proved the first
two. He died the next year, aged thirty-two.

Why should partitions, which are about adding, know anything about divisibility, which is about multiplying? The answer runs
through generating functions: $\sum p(n)x^n=\prod_k(1-x^k)^{-1}$, and a product has arithmetic hidden in it. In 1944 a young
Freeman Dyson guessed a combinatorial explanation, a statistic he called the \emph{rank}, that splits the partitions of $5n+4$
into five classes of equal size. It was proved by Atkin and Swinnerton-Dyer ten years later, and the statistic that does the
same for $11$, the \emph{crank}, was found by Andrews and Garvan only in 1988. Problem~\ref{pb:ramanujan5} proves the congruence
modulo $5$ in Ramanujan's way, with nothing more than two famous identities for infinite products.

This chapter is about such meetings. A power series is a sequence in disguise; reducing it modulo a prime, factoring it,
or substituting $x^2$ for $x$ reveals arithmetic properties of the sequence that are invisible term by term.

\begin{think}
Here are $p(0),\dots,p(24)$: 1, 1, 2, 3, 5, 7, 11, 15, 22, 30, 42, 56, 77, 101, 135, 176, 231, 297, 385, 490, 627, 792,
1002, 1255, 1575. Circle the multiples of $5$, then the multiples of $7$. Which pattern would you have bet on? And is $p(n)$
more often odd or even?
\end{think}

\section{Partitions}

A \emph{partition} of $n$ is a non-increasing sequence of positive integers (its \emph{parts}) with sum $n$; $p(n)$ is their number,
and $p(0)=1$. Its Ferrers diagram has one row of dots per part, and reflecting it in the diagonal gives the \emph{conjugate}
partition (Volume~4, Chapter~2). The basic identities are Euler's: $\sum_np(n)x^n=\prod_{k\ge1}(1-x^k)^{-1}$, and the pentagonal
number theorem $\prod_{k\ge1}(1-x^k)=\sum_{j\in\Z}(-1)^jx^{j(3j-1)/2}$ (Volume~4, Proposition~2.11 and Theorem~2.12). We write
$E(x)=\prod_{k\ge1}(1-x^k)$.

\begin{exercise}[$\star\star$]\label{pb:parts123}
\emph{Parts $1$, $2$, $3$.} Let $a(n)$ be the number of partitions of $n$ into parts equal to $1$, $2$ or $3$.
\begin{enumerate}[label=(\alph*)]
\item Show that $a(n)-a(n-6)=n$ for $n\ge6$.
\item Deduce that $a(n)$ is the integer nearest to $\frac{(n+3)^2}{12}$ for every $n\ge0$.
\item Show that the number of triangles with integer sides and perimeter $2m+3$, up to congruence, is $a(m)$, and that for
perimeter $2m$ it is $a(m-3)$.
\end{enumerate}
\end{exercise}

\begin{exercise}[$\star$]\label{pb:selfconj}
\emph{Parity by reflection.} A partition is \emph{self-conjugate} if it equals its conjugate. Let $sc(n)$ be the number of
self-conjugate partitions of $n$.
\begin{enumerate}[label=(\alph*)]
\item Show that $p(n)\equiv sc(n)\pmod2$.
\item Show that $sc(n)$ is the number of partitions of $n$ into distinct odd parts, so that
$\sum_np(n)x^n\equiv\prod_{k\ge1}(1+x^{2k-1})\pmod2$.
\item Show that
\[
\prod_{k\ge1}(1+x^{2k-1})=\sum_{k\ge0}\frac{x^{k^2}}{(1-x^2)(1-x^4)\cdots(1-x^{2k})}.
\]
\end{enumerate}
\end{exercise}

\begin{exercise}[$\star\star$]\label{pb:glaisherthm}
\emph{Glaisher's theorem.} Let $k\ge2$.
\begin{enumerate}[label=(\alph*)]
\item Show that the partitions of $n$ with no part divisible by $k$ are as many as the partitions of $n$ in which no part
appears $k$ times or more.
\item Find an explicit bijection between the two sets, and apply it to the partition $7+2+2+2+2+1+1+1$ for $k=3$.
\end{enumerate}
\end{exercise}

\begin{exercise}[$\star\star$]\label{pb:stanley}
\emph{Counting ones.}
\begin{enumerate}[label=(\alph*)]
\item (Stanley.) Show that the total number of parts equal to $1$, over all partitions of $n$, equals the total number of
\emph{distinct} part sizes, over all partitions of $n$. Check it for $n=4$.
\item (Elder.) More generally, for $k\ge1$, show that the total number of parts equal to $k$ over all partitions of $n$ equals the
total number of part sizes that appear at least $k$ times, over all partitions of $n$.
\end{enumerate}
\end{exercise}

\begin{exercise}[$\star\star$]\label{pb:eulerlog}
\emph{The logarithmic derivative.} Let $(a_d)_{d\ge1}$ be complex numbers and define $(b_n)$ by
\[
\sum_{n\ge0}b_nx^n=\prod_{d\ge1}(1-x^d)^{-a_d}.
\]
\begin{enumerate}[label=(\alph*)]
\item Show that $nb_n=\sum_{k=1}^nc_kb_{n-k}$ for $n\ge1$, where $c_k=\sum_{d\mid k}d\,a_d$.
\item Deduce that $np(n)=\sum_{k=1}^n\sigma(k)p(n-k)$, where $\sigma(k)$ is the sum of the divisors of $k$. Check it for $n=4$.
\item Let $M(d,q)$ be the number of monic irreducible polynomials of degree $d$ over $\mathbb F_q$. Using unique factorisation in
$\mathbb F_q[X]$, show that $\prod_d(1-x^d)^{-M(d,q)}=\frac1{1-qx}$, and deduce Gauss's formula $\sum_{d\mid k}d\,M(d,q)=q^k$.
\end{enumerate}
\end{exercise}

\begin{exercise}[$\star\star$]\label{pb:kolberg}
\emph{Both parities, forever (Kolberg, 1959).} Show that $p(n)$ is even for infinitely many $n$ and odd for infinitely many $n$.
\end{exercise}

\begin{figure}[ht]
\centering
\fig{ch07-ramanujan}
\caption{The residues of $p(5m+r)$ modulo $5$ for $0\le m<24$, one row for each $r$; white means $0$. The row $r=4$ is blank: this is
Ramanujan's congruence $p(5m+4)\equiv0\pmod5$ (Problem~\ref{pb:ramanujan5}). The other rows show no visible pattern.}
\label{fig:ramanujan}
\end{figure}

\begin{exercise}[Journey]\label{pb:ramanujan5}
\emph{Ramanujan's congruence.} We prove that $p(5n+4)\equiv0\pmod5$ for every $n\ge0$. Fix $q$ with $0<|q|<1$ and let, for
$z\in\C^*$,
\[
H(z)=\prod_{n\ge1}(1+zq^{2n-1})(1+z^{-1}q^{2n-1})=\sum_{m\in\Z}c_mz^m
\]
(a Laurent series, the coefficients depending on $q$).
\begin{enumerate}[label=(\alph*)]
\item Show that $H(q^2z)=\frac{H(z)}{qz}$, and deduce that $c_m=q^{m^2}c_0$ for every $m\in\Z$.
\item By evaluating at $z=i$, and at $z=-1$ with $q$ replaced by $q^4$, show that $c_0(q)E(q^2)=c_0(q^4)E(q^8)$, and deduce the
\emph{Jacobi triple product}
\[
\prod_{n\ge1}(1-q^{2n})(1+zq^{2n-1})(1+z^{-1}q^{2n-1})=\sum_{m\in\Z}q^{m^2}z^m.
\]
\item Deduce \emph{Jacobi's identity} $E(x)^3=\sum_{k\ge0}(-1)^k(2k+1)x^{k(k+1)/2}$ for $|x|<1$.
\item Show that in $xE(x)^4=xE(x)\cdot E(x)^3$, every coefficient of a power $x^{5m}$ is divisible by $5$.
\item Using $E(x)^5\equiv E(x^5)\pmod5$, conclude.
\end{enumerate}
\end{exercise}

\section{Binomial coefficients modulo primes}

Lucas's theorem (Volume~1, Exercise~3.9) reads $\binom nk$ modulo $p$ off the base-$p$ digits of $n$ and $k$. The problems of this
section push the idea further: to the exact power of $p$, to residues modulo $3$ and $4$, and to sequences built from
binomial coefficients, where a generating function does the bookkeeping. We write $s_p(n)$ for the sum of the base-$p$ digits
of $n$ and $s(n)=s_2(n)$.

\begin{exercise}[$\star\star$]\label{pb:kummer}
\emph{Kummer's theorem (1852).} Let $p$ be prime.
\begin{enumerate}[label=(\alph*)]
\item Show that $v_p\binom{m+n}m$ is the number of carries when $m$ and $n$ are added in base $p$.
\item Deduce that every prime power dividing $\binom nk$ is at most $n$.
\item Deduce that $\binom nk$ divides $\operatorname{lcm}(1,2,\dots,n)$, and that $\operatorname{lcm}(1,\dots,n)\ge\frac{2^n}{n+1}$.
\end{enumerate}
\end{exercise}

\begin{exercise}[$\star\star$]\label{pb:catalanodd}
\emph{Odd Catalan numbers.} Let $C_n=\frac1{n+1}\binom{2n}n$ be the Catalan numbers.
\begin{enumerate}[label=(\alph*)]
\item Show that $v_2(C_n)=s(n+1)-1$. Deduce that $C_n$ is odd if and only if $n=2^k-1$ for some $k\ge0$.
\item Give a second proof of the last statement from the equation $C(x)=1+xC(x)^2$ satisfied by $C(x)=\sum C_nx^n$, reduced
modulo $2$.
\end{enumerate}
\end{exercise}

\begin{figure}[ht]
\centering
\fig{ch07-pascal3}
\caption{Pascal's triangle modulo $3$, rows $0$ to $26$: blue dots are entries $\equiv1$, red dots entries $\equiv2$, and the gaps
are multiples of $3$. Row $n$ has $2^{a_1}3^{a_2}$ dots, where $a_j$ counts the digits $j$ of $n$ in base $3$, and the blue dots
always win, by exactly $2^{a_1}$ (Problem~\ref{pb:pascalthree}).}
\label{fig:pascal3}
\end{figure}

\begin{exercise}[$\star\star$]\label{pb:pascalthree}
\emph{Pascal's triangle modulo $3$.}
\begin{enumerate}[label=(\alph*)]
\item Let $p$ be prime and $n=\sum_in_ip^i$ in base $p$. Show that exactly $\prod_i(n_i+1)$ entries of row $n$ of Pascal's triangle
are not divisible by $p$.
\item For $p=3$, let $a_j$ be the number of digits of $n$ equal to $j$. Show that row $n$ has $2^{a_1}$ more entries $\equiv1$ than
entries $\equiv2\pmod3$.
\item In the first $3^k$ rows, find the numbers of entries $\equiv0$, $\equiv1$ and $\equiv2\pmod3$, and their proportions as
$k\to\infty$.
\end{enumerate}
\end{exercise}

\begin{exercise}[$\star\star$]\label{pb:pascalfour}
\emph{Modulo $4$.} Let $n\ge1$, and let $B(n)$ be the number of occurrences of the block $10$ in the binary expansion of $n$
(for instance $B(10)=B(1010_2)=2$).
\begin{enumerate}[label=(\alph*)]
\item Show that exactly $B(n)\,2^{s(n)-1}$ entries of row $n$ are $\equiv2\pmod4$.
\item Which rows contain no entry $\equiv2\pmod4$?
\end{enumerate}
\end{exercise}

\begin{exercise}[$\star\star$]\label{pb:centralp}
\emph{Central binomial coefficients modulo $p$.} Let $p$ be an odd prime and $f(x)=\sum_{n\ge0}\binom{2n}nx^n=(1-4x)^{-1/2}$.
\begin{enumerate}[label=(\alph*)]
\item Show that $f(x)\equiv(1-4x)^{(p-1)/2}f(x^p)\pmod p$, coefficient by coefficient.
\item Deduce that $\binom{2n}n\equiv\prod_i\binom{2n_i}{n_i}\pmod p$, where $n=\sum_in_ip^i$ in base $p$.
\item Show that $p\nmid\binom{2n}n$ if and only if every base-$p$ digit of $n$ is at most $\frac{p-1}2$, and that the proportion of
such $n$ below $p^k$ tends to $0$.
\end{enumerate}
\end{exercise}

\begin{exercise}[$\star\star$]\label{pb:touchard}
\emph{Touchard's congruence.} The Bell number $B_n$ is the number of partitions of the set $\{1,\dots,n\}$ into non-empty blocks.
Let $p$ be prime. By letting the cyclic group of order $p$ rotate the last $p$ elements of $\{1,\dots,n+p\}$, show that
\[
B_{n+p}\equiv B_n+B_{n+1}\pmod p.
\]
\end{exercise}

\section{Sequences written in binary}

Some sequences are defined by what happens to $2n$ and $2n+1$. Their generating functions satisfy equations relating $F(x)$ to
$F(x^2)$, and modulo $2$, where $F(x)^2=F(x^2)$, such equations become algebraic.

\begin{exercise}[$\star\star$]\label{pb:stern}
\emph{Stern's diatomic sequence (1858).} Let $s(0)=0$, $s(1)=1$, $s(2n)=s(n)$ and $s(2n+1)=s(n)+s(n+1)$.
\begin{enumerate}[label=(\alph*)]
\item A \emph{hyperbinary representation} of $n$ writes $n=\sum_ie_i2^i$ with every $e_i\in\{0,1,2\}$. Show that $n$ has exactly
$s(n+1)$ of them.
\item Deduce that $\sum_{n\ge0}s(n+1)x^n=\prod_{k\ge0}\bigl(1+x^{2^k}+x^{2^{k+1}}\bigr)$.
\item Show that $\sum_{2^k\le n<2^{k+1}}s(n)=3^k$.
\item Show that $s(n)$ and $s(n+1)$ are coprime for every $n\ge0$.
\end{enumerate}
\end{exercise}

\begin{figure}[ht]
\centering
\fig{ch07-stern}
\caption{Stern's sequence $s(n)$ for $1\le n\le256$. On each stretch $[2^k,2^{k+1}]$ the maximum is the Fibonacci number $F_{k+2}$
(red dots, Problem~\ref{pb:sternfib}), reached at two points placed symmetrically, near $\frac43\cdot2^k$ and $\frac53\cdot2^k$.}
\label{fig:stern}
\end{figure}

\begin{exercise}[$\star\star$]\label{pb:sternfib}
\emph{Stern meets Pascal and Fibonacci.} Keep the notation of Problem~\ref{pb:stern}, and let $F_n$ be the Fibonacci numbers.
\begin{enumerate}[label=(\alph*)]
\item Show that $s(n+1)$ is the number of $k$ such that $\binom{n-k}k$ is odd.
\item Show that $s(n)\le F_{k+2}$ for $2^k\le n\le2^{k+1}$.
\item Show that equality holds for $n=\frac13\bigl(2^{k+2}-(-1)^k\bigr)$.
\end{enumerate}
\end{exercise}

\begin{exercise}[$\star\star$]\label{pb:binarypart}
\emph{Binary partitions (Euler).} Let $b(n)$ be the number of partitions of $n$ into powers of $2$, so $b(0)=b(1)=1$, $b(2)=2$.
\begin{enumerate}[label=(\alph*)]
\item Show that $b(2n+1)=b(2n)$ and $b(2n)=b(2n-1)+b(n)$ for $n\ge1$.
\item Show in two ways that $b(n)$ is even for every $n\ge2$: from (a), and from the generating function, reduced modulo $2$.
\item Show that $b(2n)=b(0)+b(1)+\dots+b(n)$, and deduce that $\log b(n)\le\frac{(\log n)^2}{2\log2}+O(\log n)$.
\end{enumerate}
\end{exercise}

\begin{exercise}[$\star$]\label{pb:fibdivides}
\emph{Every integer divides a Fibonacci number.}
\begin{enumerate}[label=(\alph*)]
\item Let $m\ge1$. Show that the sequence $(F_n\bmod m)$ is periodic from the start, and deduce that $m$ divides some $F_n$ with
$n\ge1$.
\item Show that the set of $n$ with $m\mid F_n$ is the set of multiples of the least such $n\ge1$.
\item Let $J_0=0$, $J_1=1$, $J_{n+2}=J_{n+1}+2J_n$. Show that $J_n$ is odd for every $n\ge1$. Which step of (a) fails?
\end{enumerate}
\end{exercise}

\begin{exercise}[$\star\star$]\label{pb:thuemorse}
\emph{The Thue--Morse series is algebraic.} Let $t_n=s(n)\bmod2\in\mathbb F_2$ (here $s(n)$ is again the binary digit sum) and
$T(x)=\sum_nt_nx^n\in\mathbb F_2[[x]]$.
\begin{enumerate}[label=(\alph*)]
\item Show that $(1+x)^3T^2+(1+x)^2T+x=0$ in $\mathbb F_2[[x]]$.
\item Show that $(t_n)$ is not eventually periodic.
\item Deduce that $T$ is algebraic of degree $2$ over $\mathbb F_2(x)$ but not rational.
\end{enumerate}
\end{exercise}

\section{Dirichlet series}

A Dirichlet series $\sum a_nn^{-s}$ is the multiplicative cousin of a power series: products of series correspond to the
convolution $(a*b)(n)=\sum_{d\mid n}a(d)b(n/d)$, and an Euler product over the primes records a multiplicative function
(Volume~1, Chapter~4; Volume~5, Chapter~2). The four problems of this section read averages off such products.

\begin{exercise}[$\star\star$]\label{pb:orderedfact}
\emph{Ordered factorisations (Kalmár, 1931).} Let $H(1)=1$ and, for $n\ge2$, let $H(n)$ be the number of ways of writing $n$ as an
ordered product of factors $>1$ (so $H(12)=8$). Let $\rho>1$ be the solution of $\zeta(\rho)=2$; numerically $\rho=1.7286\ldots$
\begin{enumerate}[label=(\alph*)]
\item Show that $\sum_nH(n)n^{-s}=\frac1{2-\zeta(s)}$ for real $s>\rho$.
\item Compute $H(p^a)$ for a prime $p$, and $H(p_1\cdots p_k)$ for distinct primes, for $k\le4$.
\item Show that $H(n)\le n^\rho$ for every $n\ge1$.
\item Show that $\sum_{n\le x}H(n)\le\frac{x^\sigma}{2-\zeta(\sigma)}$ for every $\sigma>\rho$ and $x\ge1$.
\end{enumerate}
\end{exercise}

\begin{exercise}[$\star\star$]\label{pb:dsquared}
\emph{Ramanujan's $\sum d(n)^2$.} Let $d(n)$ be the number of divisors of $n$.
\begin{enumerate}[label=(\alph*)]
\item Show that $\sum_{n\ge1}d(n)^2n^{-s}=\frac{\zeta(s)^4}{\zeta(2s)}$ for $s>1$.
\item Deduce that $\sum_{n\le x}d(n)^2=\frac1{\pi^2}x\log^3x+O(x\log^2x)$.
\end{enumerate}
\end{exercise}

\begin{exercise}[$\star\star$]\label{pb:powerful}
\emph{Powerful numbers.} An integer $n\ge1$ is \emph{powerful} if $p^2\mid n$ for every prime $p\mid n$.
\begin{enumerate}[label=(\alph*)]
\item Show that every powerful $n$ can be written uniquely as $a^2b^3$ with $b$ squarefree.
\item Show that the number of powerful $n\le x$ is $\frac{\zeta(3/2)}{\zeta(3)}\sqrt x+O(x^{1/3})$, and that their Dirichlet series is
$\frac{\zeta(2s)\zeta(3s)}{\zeta(6s)}$.
\item Show that there are infinitely many pairs of consecutive powerful numbers, such as $(8,9)$ and $(288,289)$.
\end{enumerate}
\end{exercise}

\begin{exercise}[$\star\star$]\label{pb:abelian}
\emph{How many abelian groups? (Erdős and Szekeres, 1934).} Let $A(n)$ be the number of abelian groups of order $n$, up to
isomorphism.
\begin{enumerate}[label=(\alph*)]
\item Show that $A(n)=\prod_ip(e_i)$ if $n=\prod_ip_i^{e_i}$, and that $\sum_nA(n)n^{-s}=\prod_{k\ge1}\zeta(ks)$ for $s>1$.
\item Show that $\sum_{n\le x}A(n)=cx+O(\sqrt x)$ with $c=\prod_{k\ge2}\zeta(k)=2.2948\ldots$: on average, there are about $2.29$
abelian groups of each order.
\end{enumerate}
\end{exercise}

\section{Hints}

\begin{hint}{pb:parts123}
(a) $(1-x^6)=(1-x^3)(1+x^3)$, and $\frac1{(1-x)(1-x^2)}$ counts partitions into $1$'s and $2$'s. (b) Both sides change by $n$
when $n$ increases by $6$. (c) For a triangle with sides $x,y,z$, look at $y+z-x$, $x+z-y$, $x+y-z$.
\end{hint}

\begin{hint}{pb:selfconj}
(a) Conjugation is an involution. (b) Peel a self-conjugate diagram into hooks around the diagonal. (c) Classify by the size of
the largest square inside the diagram.
\end{hint}

\begin{hint}{pb:glaisherthm}
(a) $\frac{1-x^{kj}}{1-x^j}=1+x^j+\dots+x^{(k-1)j}$. (b) Write the multiplicity of each part in base $k$.
\end{hint}

\begin{hint}{pb:stanley}
Count the partitions of $n$ that contain at least $j$ copies of the part $m$.
\end{hint}

\begin{hint}{pb:eulerlog}
(a) Apply $x\frac d{dx}\log$ to both sides. (c) A monic polynomial is a product of monic irreducibles in exactly one way.
\end{hint}

\begin{hint}{pb:kolberg}
Otherwise $\sum p(n)x^n$ would be rational modulo $2$, and so would $E(x)$. What does a rational series over $\mathbb F_2$ look
like?
\end{hint}

\begin{hint}{pb:ramanujan5}
(b) At $z=i$ the odd powers cancel. (c) Put $q=x^{1/2}$, $z=-wx^{1/2}$, divide by $1-w^{-1}$ and let $w\to1$. (d) Multiply the
exponent by $8$ and complete the squares modulo $5$.
\end{hint}

\begin{hint}{pb:kummer}
(a) Legendre: $v_p(n!)=\frac{n-s_p(n)}{p-1}$. (b) A carry out of position $i$ needs a digit of $n$ in position $i+1$.
\end{hint}

\begin{hint}{pb:catalanodd}
(a) $s(n+1)=s(n)+1-v_2(n+1)$. (b) Modulo $2$, $C(x)^2=C(x^2)$.
\end{hint}

\begin{hint}{pb:pascalthree}
(b) Lucas, and the sign $\varepsilon(1)=1$, $\varepsilon(2)=-1$, which is multiplicative on $(\Z/3\Z)^*$.
\end{hint}

\begin{hint}{pb:pascalfour}
By Kummer, $\binom nk\equiv2\pmod4$ means exactly one carry in $k+(n-k)$. Where can it happen?
\end{hint}

\begin{hint}{pb:centralp}
(a) $f^2(1-4x)=1$ and $f(x)^p\equiv f(x^p)$. (b) $\binom{-1/2}j(-4)^j=\binom{2j}j$.
\end{hint}

\begin{hint}{pb:touchard}
In a partition fixed by the rotation, the $p$ rotated elements are either all in one block or all singletons.
\end{hint}

\begin{hint}{pb:stern}
(a) Look at $e_0$. (c) Induction on $k$. (d) Euclid.
\end{hint}

\begin{hint}{pb:sternfib}
(a) $\binom{n-k}k$ is odd if and only if $k$ and $n-2k$ have no binary digit $1$ in common. (b) Of two consecutive integers, one is
even.
\end{hint}

\begin{hint}{pb:binarypart}
(a) Does the partition contain a $1$? (b) $(1-x)B(x)=B(x^2)$. (c) $b$ is non-decreasing.
\end{hint}

\begin{hint}{pb:fibdivides}
(a) The map $(u,v)\mapsto(v,u+v)$ on $(\Z/m\Z)^2$ is a bijection. (b) $\gcd(F_a,F_b)=F_{\gcd(a,b)}$.
\end{hint}

\begin{hint}{pb:thuemorse}
(a) $t_{2n}=t_n$, $t_{2n+1}=1+t_n$, and $T(x)^2=T(x^2)$. (b) A period $P$ could be halved if even; if it is odd, compare
$t_{2n+P}$ and $t_{2n+1+P}$.
\end{hint}

\begin{hint}{pb:orderedfact}
(a) $H=\sum_k(\zeta-1)^k$. (c) Strong induction, with $H(n)=\sum_{d\mid n,\,d<n}H(d)$. (d) Rankin's trick.
\end{hint}

\begin{hint}{pb:dsquared}
(a) Compare Euler factors. (b) $d^2=\tau_4*g$ with $g$ supported on squares, and Volume~5, Exercise~2.7.
\end{hint}

\begin{hint}{pb:powerful}
(a) Even exponents go into $a^2$. (c) Pell's equation $x^2-8y^2=1$.
\end{hint}

\begin{hint}{pb:abelian}
(a) Volume~2, Corollary~7.7. (b) $A=\mathbf1*g$ with $g\ge0$ supported on powerful numbers.
\end{hint}

\section{Solutions}

\begin{solution}{pb:parts123}
\idea The generating function $\frac1{(1-x)(1-x^2)(1-x^3)}$ has a pole of order $3$ at $1$, which gives the quadratic
$\frac{(n+3)^2}{12}$, and smaller poles at roots of unity of order $2$, $3$, $6$, whose contributions are periodic and small. A
difference of step $6$ kills them all.

\writeup (a) The generating function of $a$ is $A(x)=\frac1{(1-x)(1-x^2)(1-x^3)}$, so
\[
(1-x^6)A(x)=\frac{1+x^3}{(1-x)(1-x^2)}.
\]
The coefficient of $x^n$ in $\frac1{(1-x)(1-x^2)}$ is the number of ways to write $n=u+2v$, which is $\lfloor\frac n2\rfloor+1$. So for
$n\ge6$,
\[
a(n)-a(n-6)=\Bigl\lfloor\frac n2\Bigr\rfloor+1+\Bigl\lfloor\frac{n-3}2\Bigr\rfloor+1=n,
\]
as one checks separately for $n$ even ($\frac n2+1+\frac{n-4}2+1$) and odd ($\frac{n-1}2+1+\frac{n-3}2+1$).

(b) Let $g(n)=\frac{(n+3)^2}{12}$. Then $g(n)-g(n-6)=\frac{(n+3)^2-(n-3)^2}{12}=n$, so $a(n)-g(n)$ depends only on $n\bmod6$. For
$n=0,\dots,5$ we have $a(n)=1,1,2,3,4,5$ and $g(n)=\frac34,\frac43,\frac{25}{12},3,\frac{49}{12},\frac{16}3$, so
$a(n)-g(n)\in\{\frac14,-\frac13,-\frac1{12},0,-\frac1{12},-\frac13\}$. These are all less than $\frac12$ in absolute value, so $a(n)$ is
the integer nearest to $g(n)$, for every $n$.

(c) Let $x\le y\le z$ be the sides, $x+y+z=n$, $x+y>z$. Put $\alpha=y+z-x\ge\beta=x+z-y\ge\gamma=x+y-z\ge1$. These are
positive, have the parity of $n$, add up to $n$, and determine $(x,y,z)=(\frac{\beta+\gamma}2,\frac{\alpha+\gamma}2,\frac{\alpha+\beta}2)$;
conversely any such triple gives a triangle. If $n=2m+3$, write $\alpha=2a'+1$, $\beta=2b'+1$, $\gamma=2c'+1$: then
$a'\ge b'\ge c'\ge0$ and $a'+b'+c'=m$, a partition of $m$ into at most $3$ parts. By conjugation these are as many as the
partitions of $m$ into parts at most $3$, which are counted by $a(m)$. If $n=2m$, write $\alpha=2a''$, $\beta=2b''$, $\gamma=2c''$
with $a''\ge b''\ge c''\ge1$ and sum $m$; subtracting $1$ from each gives a partition of $m-3$ into at most $3$ parts.

\tools Volume~4, Chapter~2 (generating functions, partial fractions, conjugation).
\leads The sequence of (c) is known as Alcuin's sequence: it also counts the solutions of a cask-sharing puzzle in the
\emph{Propositiones ad acuendos juvenes}, a collection attributed to Alcuin of York (OEIS A005044). For parts $1,\dots,k$ the
same method gives $a_k(n)=\frac{n^{k-1}}{k!(k-1)!}+O(n^{k-2})$, but the nearest-integer formula is special to small $k$.
\end{solution}

\begin{solution}{pb:selfconj}
\idea An involution changes the parity of a finite set only through its fixed points. The fixed points of conjugation have a
rigid shape that is easy to count.

\writeup (a) Conjugation is an involution of the set of partitions of $n$. The partitions that are not fixed come in pairs
$\{\lambda,\lambda'\}$, so $p(n)$ and the number $sc(n)$ of fixed points have the same parity.

(b) Let $\lambda$ be self-conjugate, and let $k$ be the number of dots on the diagonal of its Ferrers diagram. The $i$-th
\emph{diagonal hook} consists of the dot $(i,i)$, the dots to its right in row $i$, and the dots below it in column $i$. Since the
diagram is symmetric, the arm and the leg of this hook have the same length $\lambda_i-i$, so the hook has $2(\lambda_i-i)+1$ dots, an
odd number, and these numbers strictly decrease with $i$. The hooks partition the diagram, so we get a partition of $n$ into
distinct odd parts. Conversely, distinct odd parts $h_1>\dots>h_k$ can be bent into nested symmetric hooks, each fitting outside
the previous one because $h_i\ge h_{i+1}+2$; this is the inverse map. The generating function of partitions into distinct odd
parts is $\prod_k(1+x^{2k-1})$.

(c) A self-conjugate partition whose largest inscribed square (the Durfee square) is $k\times k$ consists of that square,
of the partition $\mu$ formed by the rows below the square, whose parts are at most $k$, and of its mirror image $\mu'$ to the
right of the square. So it is determined by $k$ and by a partition $\mu$ into parts at most $k$, and its size is $k^2+2|\mu|$. Summing over $k$,
the generating function of self-conjugate partitions is $\sum_kx^{k^2}\prod_{j\le k}(1-x^{2j})^{-1}$. Compare with (b).

\tools Volume~4, Chapter~2 (Ferrers diagrams, conjugation).
\leads Parity of $p(n)$ is a famous mystery: it is conjectured that $p(n)$ is odd for asymptotically half of all $n$, but it is
not even known that the odd values have a positive lower density. Problem~\ref{pb:kolberg} shows at least that both parities occur
infinitely often.
\end{solution}

\begin{solution}{pb:glaisherthm}
\idea Generating functions give the equality of the counts in one line. For the bijection, a part repeated $m$ times can be
merged into bigger parts according to the base-$k$ digits of $m$.

\writeup (a) The generating function of partitions with no part divisible by $k$ is $\prod_{k\nmid j}(1-x^j)^{-1}$; that of
partitions with every multiplicity at most $k-1$ is
\[
\prod_{j\ge1}(1+x^j+\dots+x^{(k-1)j})=\prod_{j\ge1}\frac{1-x^{kj}}{1-x^j}.
\]
In the last product, every factor $1-x^{kj}$ of the numerator cancels the factor of the denominator with index $kj$, and what is
left is $\prod_{k\nmid j}(1-x^j)^{-1}$. The two products are equal in $\Z[[x]]$, so the counts agree for every $n$.

(b) Start from a partition with no part divisible by $k$. If the part $j$ appears $m$ times, write $m=\sum_ia_ik^i$ in base $k$, and
replace the $m$ copies of $j$ by $a_i$ copies of $jk^i$ for each $i$; the sum is unchanged. Every integer is uniquely $jk^i$ with
$k\nmid j$, so the new parts $jk^i$ are distinct for distinct pairs $(j,i)$, and each appears $a_i\le k-1$ times. Conversely, from a
partition with multiplicities at most $k-1$, write each part as $jk^i$ with $k\nmid j$, and break each copy of $jk^i$ into $k^i$
copies of $j$; the base-$k$ expansion of the resulting multiplicity of $j$ recovers the digits. The two maps are inverse.

For $7+2+2+2+2+1+1+1$ (no part divisible by $3$): the part $7$ appears once, $1=1_3$, and stays; the part $2$ appears $4=11_3$
times and becomes $2+6$; the part $1$ appears $3=10_3$ times and becomes $3$. The image is $7+6+3+2$, a partition of $18$ with all
multiplicities at most $2$, as is the original.

\tools Volume~4, Chapter~2 (generating functions of partitions); Volume~4, Proposition~2.11 (the case $k=2$).
\leads For $k=2$ this is Euler's ``distinct parts equal odd parts'', and (b) is its binary bijection; Sylvester found a very
different one. Glaisher published the general case in 1883. Many partition identities of this kind, such as the Rogers--Ramanujan identities, have generating
function proofs long before any bijection was found.
\end{solution}

\begin{solution}{pb:stanley}
\idea Count the same pairs in two ways: a partition together with ``a $j$-th copy of the part $m$''.

\writeup (b) contains (a) as the case $k=1$, so we prove (b). Fix $m\ge1$ and $j\ge1$. Removing $j$ copies of $m$ is a bijection
from the partitions of $n$ with at least $j$ copies of $m$ onto the partitions of $n-jm$, so there are $p(n-jm)$ of them (with
$p$ of a negative number equal to $0$).

The total number of parts equal to $k$ over all partitions of $n$ is $\sum_\lambda m_k(\lambda)$, where $m_k$ is the multiplicity of $k$.
Since $m_k(\lambda)=\#\{j\ge1:m_k(\lambda)\ge j\}$, this total is $\sum_{j\ge1}\#\{\lambda:m_k(\lambda)\ge j\}=\sum_{j\ge1}p(n-jk)$.

The total number of part sizes appearing at least $k$ times is $\sum_{m\ge1}\#\{\lambda:m_m(\lambda)\ge k\}=\sum_{m\ge1}p(n-km)$. The two
sums are the same sum, with the summation variable renamed.

For $n=4$ the partitions are $4$, $3+1$, $2+2$, $2+1+1$, $1+1+1+1$. They contain $0+1+0+2+4=7$ ones and $1+2+1+2+1=7$
distinct part sizes.

\tools Volume~0, Chapter~8 (double counting).
\leads (a) was found by Richard Stanley, and (b) by Paul Elder in 1984; Honsberger's \emph{Mathematical Gems~III} tells the
story. The common value in (a) is $\sum_{j\ge1}p(n-j)$, so the average number of $1$'s in a partition of $n$ is
$\frac{p(0)+\dots+p(n-1)}{p(n)}$, which grows like $\frac{\sqrt{6n}}\pi$ by the Hardy--Ramanujan asymptotics.
\end{solution}

\begin{solution}{pb:eulerlog}
\idea The logarithm turns the infinite product into a sum, and $x\frac d{dx}$ turns it back into a power series with integer
structure. Multiplying back by $B(x)$ gives a recurrence.

\writeup (a) Let $B(x)=\sum b_nx^n$. As formal power series with constant term $1$,
\[
\log B(x)=\sum_da_d\log\frac1{1-x^d}=\sum_da_d\sum_{r\ge1}\frac{x^{dr}}r,\qquad x\frac{B'(x)}{B(x)}=\sum_da_d\sum_{r\ge1}dx^{dr}=\sum_{k\ge1}c_kx^k,
\]
with $c_k=\sum_{d\mid k}da_d$ (put $k=dr$). So $xB'(x)=B(x)\sum_kc_kx^k$, and comparing coefficients of $x^n$ gives
$nb_n=\sum_{k=1}^nc_kb_{n-k}$. (One can avoid logarithms: the logarithmic derivative of a product is the sum of the logarithmic
derivatives of the factors, and $x\frac d{dx}$ of $(1-x^d)^{-a_d}$ divided by itself is $\frac{a_ddx^d}{1-x^d}$.)

(b) Take $a_d=1$ for all $d$: then $b_n=p(n)$ and $c_k=\sigma(k)$. For $n=4$: $\sigma(1)p(3)+\sigma(2)p(2)+\sigma(3)p(1)+\sigma(4)p(0)
=3+3\cdot2+4+7=20=4\,p(4)$.

(c) Expand $\prod_P(1-x^{\deg P})^{-1}$ over all monic irreducible $P$: a term $x^n$ arises from each choice of finitely many $P$ with
multiplicities and total degree $n$, that is, by unique factorisation, from each monic polynomial of degree $n$. There are $q^n$ of
them, so the product is $\sum q^nx^n=\frac1{1-qx}$, and grouping the $P$ by degree gives $\prod_d(1-x^d)^{-M(d,q)}$. Now apply
(a) with $a_d=M(d,q)$, $b_n=q^n$; or directly: $x\frac{d}{dx}\log\frac1{1-qx}=\frac{qx}{1-qx}=\sum_kq^kx^k$, so $c_k=q^k$, that is
$\sum_{d\mid k}dM(d,q)=q^k$.

\tools Volume~4, Chapter~2 (formal logarithm and derivative); Volume~2, Chapter~9 (irreducible polynomials over $\mathbb F_q$).
\leads Möbius inversion turns Gauss's formula into $M(k,q)=\frac1k\sum_{d\mid k}\mu(k/d)q^d$. The map $(a_d)\mapsto(b_n)$ is the
\emph{Euler transform}; it turns ``connected'' objects into all objects whenever objects factor uniquely into connected ones,
as integers factor into primes and polynomials into irreducibles.
\end{solution}

\begin{solution}{pb:kolberg}
\idea If $p(n)$ had eventually constant parity, the generating function of partitions would be a rational function modulo $2$, and
so would its inverse, the pentagonal series. But the pentagonal series has longer and longer gaps, which a rational series over
a finite field cannot have.

\writeup Work in $\mathbb F_2[[x]]$ and let $\bar P=\sum_n(p(n)\bmod2)x^n$. If $p(n)$ is even for all $n\ge N$, then $\bar P$ is a polynomial;
if $p(n)$ is odd for all $n\ge N$, then $\bar P$ is a polynomial plus $\frac{x^N}{1-x}$. In both cases $\bar P=\frac UV$ with polynomials
$U,V$ and $U(0)=V(0)=1$ (as $p(0)=1$). Since $P(x)E(x)=1$ in $\Z[[x]]$, the reduction $\bar E$ of $E$ is $\frac VU$, rational too.

A rational series $\frac VU=\sum e_nx^n$ over $\mathbb F_2$ with $U=1+u_1x+\dots+u_dx^d$ satisfies $e_n=u_1e_{n-1}+\dots+u_de_{n-d}$ for
$n>\max(d,\deg V)$. The vector $(e_{n-1},\dots,e_{n-d})$ takes at most $2^d$ values, so it repeats, and from then on the sequence
is periodic. An eventually periodic sequence with infinitely many non-zero terms has bounded gaps between them.

But by the pentagonal number theorem, $e_n\ne0$ exactly when $n=\frac{j(3j\mp1)}2$ for some $j\ge0$, and the gap between the
consecutive values $\frac{j(3j+1)}2$ and $\frac{(j+1)(3j+2)}2$ is $2j+1$, unbounded. This contradiction shows that neither parity
can hold from some point on.

\tools Volume~4, Theorem~2.12 (pentagonal number theorem); Volume~4, Theorem~2.5 (rational series and recurrences).
\leads The argument gives no quantitative information. Nicolas, Ruzsa and Sárközy (1998) and Ahlgren (1999) showed that each
parity occurs, up to logarithmic factors, at least $\sqrt N$ times below $N$, still far from the expected $\frac N2$. Every arithmetic
progression contains infinitely many $n$ with $p(n)$ even (Ono, 1996) and infinitely many with $p(n)$ odd (Radu, 2012).
\end{solution}

\begin{solution}{pb:ramanujan5}
\idea Two infinite products have explicit expansions: $E(x)$ (Euler) and $E(x)^3$ (Jacobi). Their product $E(x)^4$ has
coefficients we can control modulo $5$, and $E(x)^4=E(x)^5/E(x)\equiv E(x^5)P(x)$ links them to partitions.

\writeup (a) The product converges absolutely and locally uniformly on $\C^*$ since $\sum|q|^{2n-1}<\infty$, so $H$ is holomorphic on
$\C^*$ and has a Laurent expansion. Replacing $z$ by $q^2z$ shifts the factors:
\[
H(q^2z)=\prod_{n\ge1}(1+zq^{2n+1})(1+z^{-1}q^{2n-3})=H(z)\,\frac{1+z^{-1}q^{-1}}{1+zq}=\frac{H(z)}{zq}.
\]
In coefficients, $c_mq^{2m}=q^{-1}c_{m+1}$, that is $c_{m+1}=q^{2m+1}c_m$ for all $m\in\Z$. Since $(m+1)^2-m^2=2m+1$, induction upwards
and downwards from $m=0$ gives $c_m=q^{m^2}c_0$. So $H(z)=c_0(q)\theta(z,q)$ with $\theta(z,q)=\sum_mq^{m^2}z^m$.

(b) At $z=i$, $(1+iq^{2n-1})(1-iq^{2n-1})=1+q^{4n-2}$, while in $\theta(i,q)$ the terms $m$ and $-m$ cancel for odd $m$ (as
$i^m+i^{-m}=0$) and $m=2k$ gives $(-1)^kq^{4k^2}$. So
\[
\prod_{n\ge1}(1+q^{4n-2})=c_0(q)\sum_k(-1)^kq^{4k^2}.
\]
At $z=-1$ with $q^4$ in place of $q$, $\prod_n(1-q^{8n-4})^2=c_0(q^4)\sum_k(-1)^kq^{4k^2}$. The series $\sum(-1)^kq^{4k^2}=1-2q^4+\dots$ does
not vanish for small $|q|$; dividing, and using $1-q^{8n-4}=(1-q^{4n-2})(1+q^{4n-2})$,
\[
\frac{c_0(q)}{c_0(q^4)}=\frac1{\prod_n(1-q^{4n-2})\prod_n(1-q^{8n-4})}.
\]
Splitting $E(q^2)=\prod(1-q^{2n})$ according to the parity of $n$, and then $\prod(1-q^{4n})$ in the same way,
$E(q^2)=E(q^8)\prod(1-q^{8n-4})\prod(1-q^{4n-2})$. Hence $d(q)=c_0(q)E(q^2)$ satisfies $d(q)=d(q^4)$ for small $|q|$. Now
$c_0(q)=\frac1{2\pi}\int_0^{2\pi}H(e^{it})\,dt$ is holomorphic in $q$ on the unit disc and tends to $1$ as $q\to0$ (the product tends to
$1$ uniformly on $|z|=1$). So $d(q)=d(q^{4^k})\to1$ as $k\to\infty$, $d=1$ near $0$, and by analytic continuation $c_0(q)=\frac1{E(q^2)}$
on the whole disc. Multiplying $H(z)=c_0(q)\theta(z,q)$ by $E(q^2)$ gives the triple product.

(c) Take $0<x<1$, $q=x^{1/2}$ and $z=-wx^{1/2}$ with $w>0$. Then $1+zq^{2n-1}=1-wx^n$, $1+z^{-1}q^{2n-1}=1-w^{-1}x^{n-1}$ and
$q^{m^2}z^m=(-1)^mw^mx^{m(m+1)/2}$, so
\[
E(x)\prod_{n\ge1}(1-wx^n)\prod_{n\ge1}(1-w^{-1}x^{n-1})=\sum_{m\in\Z}(-1)^mw^mx^{m(m+1)/2}.
\]
The factor $n=1$ of the last product is $1-w^{-1}$. On the right, the terms $m$ and $-1-m$ have the same exponent of $x$ and
opposite signs, so the right side is $\sum_{m\ge0}(-1)^mx^{m(m+1)/2}(w^m-w^{-m-1})$. Divide both sides by $1-w^{-1}$ and let $w\to1$:
$\frac{w^m-w^{-m-1}}{1-w^{-1}}\to2m+1$ (both numerator and denominator vanish at $w=1$, with derivatives $2m+1$ and $1$), and the left
side tends to $E(x)\cdot E(x)\cdot E(x)$. The exchange of limit and sum is justified by uniform convergence for $w\in[\frac12,2]$.
Both sides of Jacobi's identity are power series in $x$ that agree on $]0,1[$, so they agree for $|x|<1$.

(d) By (c) and the pentagonal theorem,
\[
xE(x)^4=\sum_{j\in\Z}\sum_{k\ge0}(-1)^{j+k}(2k+1)x^{e(j,k)},\qquad e(j,k)=1+\frac{j(3j-1)}2+\frac{k(k+1)}2.
\]
Now
\[
8e(j,k)=(2k+1)^2+12j^2-4j+7\equiv(2k+1)^2+2(j+4)^2\pmod5,
\]
since $2(j+4)^2=2j^2+16j+32\equiv2j^2+j+2\equiv12j^2-4j+7$. The squares modulo $5$ are $0,1,4$, and twice them are $0,2,3$; a
square plus twice a square is $\equiv0$ only if both are $0$. So $5\mid e(j,k)$ forces $5\mid2k+1$, and every term contributing to a
power $x^{5m}$ has a coefficient divisible by $5$.

(e) In $\Z[[x]]$, $(1-x^k)^5\equiv1-x^{5k}\pmod5$ (the middle binomial coefficients are divisible by $5$), so $E(x)^5\equiv E(x^5)$ and
\[
xE(x)^4=\frac{xE(x)^5}{E(x)}\equiv E(x^5)\sum_{n\ge0}p(n)x^{n+1}\pmod5.
\]
Split $\sum p(n)x^{n+1}=\sum_{r=0}^4x^rA_r(x^5)$ by the residue of the exponent modulo $5$. Multiplying by $E(x^5)$ does not mix the
residues, so the part of $xE(x)^4$ with exponents divisible by $5$ is $E(x^5)A_0(x^5)$, which is $\equiv0\pmod5$ by (d). Since
$E(x^5)$ is invertible in $\Z[[x]]$, $A_0\equiv0\pmod5$. The coefficients of $A_0(x^5)$ are the $p(n)$ with $n+1\equiv0\pmod5$, that
is, the $p(5m+4)$.

\tools Volume~4, Theorem~2.12 (pentagonal number theorem); Volume~4, Chapters~2 and~9 (formal and holomorphic power series,
Laurent series); Volume~1, Chapter~3 (squares modulo $5$).
\leads This is Ramanujan's proof from 1919. The same idea with $E(x)^6=E(x)^3E(x)^3$ gives $p(7n+5)\equiv0\pmod7$, and the
congruence modulo $11$ needs more (Ramanujan's unpublished manuscript, edited by Hardy in 1921). Ramanujan also proved the
beautiful identity $\sum_np(5n+4)x^n=5\frac{E(x^5)^5}{E(x)^6}$. In 2000 Ken Ono showed that there are congruences of the form
$p(An+B)\equiv0\pmod M$ for every $M$ prime to $6$, through the theory of modular forms (Volume~5, Chapter~12, and Part~III).
\end{solution}

\begin{solution}{pb:kummer}
\idea Legendre's formula counts the factors $p$ in $n!$ through the digit sum, and each carry in an addition costs exactly $p-1$
in digit sum.

\writeup (a) By Volume~1, Exercise~12.5, $v_p(n!)=\frac{n-s_p(n)}{p-1}$. Hence
\[
v_p\binom{m+n}m=\frac{s_p(m)+s_p(n)-s_p(m+n)}{p-1}.
\]
Add $m$ and $n$ digit by digit. In a position with incoming carry $c_{i-1}\in\{0,1\}$ the digit of the sum is
$m_i+n_i+c_{i-1}-pc_i$, where $c_i$ is the outgoing carry. Summing over $i$ (with $c_{-1}=0$ and no carry beyond the last digit),
$s_p(m+n)=s_p(m)+s_p(n)+\sum_ic_{i-1}-p\sum_ic_i=s_p(m)+s_p(n)-(p-1)\sum_ic_i$. So $v_p\binom{m+n}m=\sum_ic_i$.

(b) Add $k$ and $n-k$. If $n$ has $L+1$ digits ($p^L\le n<p^{L+1}$), a carry out of position $L$ would make the sum at least
$p^{L+1}$, so carries can only occur out of positions $0,\dots,L-1$: at most $L$ of them. Hence $p^{v_p\binom nk}\le p^L\le n$.

(c) By (b), for each prime $p$ the power $p^{v_p\binom nk}$ is at most $n$, hence divides the highest power of $p$ that is $\le n$,
which is the power of $p$ in $\operatorname{lcm}(1,\dots,n)$. So $\binom nk$ divides the lcm. The largest of the $n+1$ numbers
$\binom nk$ is at least their average $\frac{2^n}{n+1}$.

\tools Volume~1, Exercise~12.5 (Legendre's formula via digit sums); Volume~1, Chapter~2 (lcm and prime powers).
\leads The bound (c) is the starting point of Chebyshev-type estimates: $\log\operatorname{lcm}(1,\dots,n)=\psi(n)$, so (c) gives
$\psi(n)\ge n\log2-\log(n+1)$ (compare Volume~5, Chapter~1). Kummer found (a) in 1852, in a long paper on reciprocity laws.
\end{solution}

\begin{solution}{pb:catalanodd}
\idea Kummer's theorem turns the parity question into a statement about binary digits; the functional equation gives the same
answer through the Frobenius map $C\mapsto C^2$.

\writeup (a) By Problem~\ref{pb:kummer}, $v_2\binom{2n}n$ is the number of carries in $n+n$ in binary, which is $s(n)$ (each digit $1$ of
$n$ produces one carry, since doubling is a shift). So $v_2(C_n)=s(n)-v_2(n+1)$. Adding $1$ to $n$ turns the final block of
$v_2(n+1)$ ones into zeros and the next zero into a one, so $s(n+1)=s(n)-v_2(n+1)+1$. Hence $v_2(C_n)=s(n+1)-1$, which vanishes
exactly when $n+1$ is a power of $2$.

(b) Reduce $C=1+xC^2$ modulo $2$. In $\mathbb F_2[[x]]$, squaring is additive, so $C(x)^2=\sum C_nx^{2n}=C(x^2)$, and
$C(x)\equiv1+xC(x^2)$. In coefficients: $C_0\equiv1$, $C_{2m}\equiv0$ for $m\ge1$, and $C_{2m+1}\equiv C_m$. By induction $C_n$ is odd if
and only if $n=0$ or $n=2m+1$ with $m=2^k-1$, that is, $n=2^{k+1}-1$.

\tools Problem~\ref{pb:kummer}; Volume~3, Exercise~1.7 (Catalan numbers); Volume~4, Chapter~2 (the equation $C=1+xC^2$).
\leads Over $\mathbb F_2$, (b) says $C(x)=\sum_kx^{2^k-1}$: a lacunary series that is a root of the quadratic $xC^2+C+1$. Such
``algebraic but sparse'' series are impossible over $\Q$: by Hadamard's gap theorem the circle of convergence of $\sum x^{2^k-1}$ is a
natural boundary, while an algebraic function has finitely many singularities. Problem~\ref{pb:thuemorse} explains what
algebraic series over $\mathbb F_p$ look like.
\end{solution}

\begin{solution}{pb:pascalthree}
\idea Lucas's theorem makes $\binom nk\bmod p$ a product of digit-by-digit factors, so any multiplicative statistic of the row
factors over the digits.

\writeup (a) By Lucas, $\binom nk\not\equiv0$ if and only if $k_i\le n_i$ for every digit, and there are $\prod_i(n_i+1)$ such $k$.

(b) Let $\varepsilon(1)=1$, $\varepsilon(2)=-1$; it is a homomorphism $(\Z/3\Z)^*\to\{\pm1\}$. Summing over the $k$ of (a),
\[
\#\{k:\tbinom nk\equiv1\}-\#\{k:\tbinom nk\equiv2\}=\sum_{k_i\le n_i}\varepsilon\Bigl(\prod_i\tbinom{n_i}{k_i}\Bigr)=\prod_i\sum_{k_i=0}^{n_i}\varepsilon\tbinom{n_i}{k_i}.
\]
The inner sum is $1$ for $n_i=0$, $1+1=2$ for $n_i=1$, and $1-1+1=1$ for $n_i=2$ (as $\binom21=2$). So the difference is $2^{a_1}$.

(c) Rows $0,\dots,3^k-1$ are the $n$ with $k$ digits $n_i\in\{0,1,2\}$. Summing the products over all choices,
the number of non-zero entries is $(1+2+3)^k=6^k$ and the difference between the $1$'s and the $2$'s is $(1+2+1)^k=4^k$. So there
are $\frac{6^k+4^k}2$ entries $\equiv1$, $\frac{6^k-4^k}2$ entries $\equiv2$, and $\frac{3^k(3^k+1)}2-6^k$ multiples of $3$. The proportion of
multiples of $3$ tends to $1$, while among the others the $1$'s and $2$'s become equally frequent, since $\frac{4^k}{6^k}\to0$.

\tools Volume~1, Exercise~3.9 (Lucas's theorem).
\leads For a general prime $p$, the multiplicative characters of $(\Z/p\Z)^*$ play the role of $\varepsilon$, and their sums over a
row give the number of entries in each residue class. This was worked out by Garfield and Wilf (1992); Granville's survey
listed below tells the story and does the same modulo prime powers.
\end{solution}

\begin{solution}{pb:pascalfour}
\idea An entry is $2\bmod4$ when it is divisible by $2$ exactly once, that is, when there is exactly one carry. One carry can only
happen at a place where $n$ has the pattern $10$.

\writeup (a) By Problem~\ref{pb:kummer}, $\binom nk\equiv2\pmod4$ if and only if adding $k$ and $m=n-k$ in binary produces exactly one carry,
say out of position $i$. In position $i$ there is no incoming carry, and $k_i+m_i=n_i+2$, which forces $k_i=m_i=1$ and $n_i=0$. In
position $i+1$ there is an incoming carry and no outgoing one, so $k_{i+1}+m_{i+1}+1=n_{i+1}\le1$, which forces $k_{i+1}=m_{i+1}=0$ and
$n_{i+1}=1$. In every other position there is no carry at all, so $k_j+m_j=n_j$: where $n_j=0$ both digits are $0$, and where $n_j=1$
we may choose which of $k_j,m_j$ is $1$. So the $k$ with exactly one carry correspond to a choice of a block $10$ in $n$ (positions
$i+1,i$) and a free choice in each of the other $s(n)-1$ positions where $n$ has a digit $1$: $B(n)2^{s(n)-1}$ of them.

(b) Rows with $B(n)=0$: binary expansions with no $1$ immediately followed by a $0$, that is, $n=11\cdots1_2=2^a-1$ (and $n=0$).

\tools Problem~\ref{pb:kummer}.
\leads There are similar rules for $\binom nk$ modulo any prime power, in terms of blocks of digits; Granville's survey listed
below explains them, and how such counts become ``fractal'' pictures as in Figure~\ref{fig:pascal3}.
\end{solution}

\begin{solution}{pb:centralp}
\idea The series $f$ satisfies $f^2(1-4x)=1$, and modulo $p$ raising to the power $p$ is the substitution $x\mapsto x^p$. This turns
the algebraic equation into a self-similarity that reads the coefficients digit by digit.

\writeup (a) For a power series $g$ with integer coefficients, $g(x)^p\equiv g(x^p)\pmod p$ (the multinomial coefficients of the
$p$-th power are divisible by $p$ except for the pure powers, and $a^p\equiv a$). Since $f^2=\frac1{1-4x}$,
\[
(1-4x)^{\frac{p-1}2}f(x^p)\equiv(1-4x)^{\frac{p-1}2}f(x)^p=(1-4x)^{\frac{p-1}2}f(x)\bigl(f(x)^2\bigr)^{\frac{p-1}2}=f(x).
\]

(b) Let
\[
Q(x)=(1-4x)^{(p-1)/2}=\sum_{j\le(p-1)/2}\binom{(p-1)/2}j(-4)^jx^j.
\]
Since $\frac{p-1}2\equiv-\frac12\pmod p$ and
$\binom{-1/2}j(-4)^j=\binom{2j}j$ (Volume~3, Exercise~3.7), the coefficient of $x^j$ in $Q$ is $\equiv\binom{2j}j$ for $j\le\frac{p-1}2$; and for
$\frac{p-1}2<j<p$, $\binom{2j}j=\frac{(2j)!}{j!^2}$ is divisible by $p$ because $p\le2j<2p$. So $Q(x)\equiv\sum_{j<p}\binom{2j}jx^j$. Writing
$n=n_0+pn'$ with $0\le n_0<p$, the coefficient of $x^n$ in $Q(x)f(x^p)$ is $\binom{2n_0}{n_0}\binom{2n'}{n'}$ modulo $p$ (only the term
$x^{n_0}$ of $Q$ can be combined with a power of $x^p$, as $\deg Q<p$). Induction on the number of digits gives the product
formula.

(c) A digit $n_i\le\frac{p-1}2$ gives $\binom{2n_i}{n_i}=\frac{(2n_i)!}{n_i!^2}$ with $2n_i<p$, not divisible by $p$; a digit $n_i>\frac{p-1}2$ gives a
multiple of $p$, as in (b). Below $p^k$ the good $n$ are $\bigl(\frac{p+1}2\bigr)^k$ out of $p^k$, a proportion
$\bigl(\frac{p+1}{2p}\bigr)^k\to0$.

\tools Volume~3, Exercise~3.7 (the series of $(1-4x)^{-1/2}$); Volume~1, Chapter~3 (Fermat's little theorem).
\leads The same trick works for any power series $f$ with $f^d\in\Q(x)$, and more generally for diagonals of rational functions.
It gives ``Lucas congruences'' for many combinatorial sequences, such as the Apéry numbers of the irrationality proof of $\zeta(3)$
(Chapter~9; Gessel, 1982). In base $p$
these series are \emph{$p$-automatic}, the theme of Problem~\ref{pb:thuemorse}.
\end{solution}

\begin{solution}{pb:touchard}
\idea A group of prime order acting on a finite set has orbits of size $1$ or $p$, so the size of the set is congruent to the
number of fixed points modulo $p$. Here the fixed partitions are easy to list.

\writeup Let $X$ be the set of partitions of $\{1,\dots,n+p\}$ and $\sigma$ the permutation that fixes $1,\dots,n$ and cycles
$n+1\to n+2\to\dots\to n+p\to n+1$. The group $\langle\sigma\rangle$ of order $p$ acts on $X$, every orbit has size $1$ or $p$, and so
$B_{n+p}=|X|\equiv\#\{\text{fixed points}\}\pmod p$.

Let $\pi$ be fixed by $\sigma$; then $\sigma$ permutes the blocks of $\pi$. Suppose some block $S$ contains some of the new elements
$n+1,\dots,n+p$ but not all. Its images $S,\sigma S,\dots,\sigma^{p-1}S$ are blocks, and they are distinct: the orbit of $S$ has size
dividing $p$, and size $1$ would mean $\sigma S=S$, so $S$ would contain the whole $\sigma$-orbit of any new element it contains,
that is, all of them. So these $p$ blocks are disjoint and all have the same intersection with $\{1,\dots,n\}$ (which $\sigma$ fixes
pointwise); being disjoint, this intersection is empty. As each $\sigma^jS$ contains a new element and there are $p$ of them,
each block $\sigma^jS$ is a single new element. So in a fixed partition, either all the new elements lie in one block, or each
new element is a singleton block.

Fixed partitions of the first kind correspond to partitions of $\{1,\dots,n,\star\}$, where $\star$ stands for the block of new
elements: $B_{n+1}$ of them. Those of the second kind correspond to partitions of $\{1,\dots,n\}$: $B_n$ of them. Both kinds are
indeed fixed. Hence $B_{n+p}\equiv B_{n+1}+B_n\pmod p$.

\tools Volume~2, Chapter~3 (orbits of a group action, orbit sizes divide the order); Volume~4, Exercise~2.3 (Bell numbers).
\leads Iterating, $B_{n+p^m}\equiv mB_n+B_{n+1}\pmod p$, and the sequence $(B_n\bmod p)$ is periodic with period dividing
$\frac{p^p-1}{p-1}$; it is conjectured that this is the exact minimal period for every prime $p$ (checked for small $p$). The same
fixed-point trick proves Fermat's little theorem with necklaces (Volume~2, Exercise~3.2).
\end{solution}

\begin{solution}{pb:stern}
\idea The recursions for $s$ and for hyperbinary representations are the same once we look at the last digit $e_0$.

\writeup (a) Let $h(n)$ be the number of hyperbinary representations of $n$, $h(0)=1$. If $n=2m+1$ is odd then $e_0=1$, and the
remaining digits form a hyperbinary representation of $m$: $h(2m+1)=h(m)$. If $n=2m\ge2$ then $e_0\in\{0,2\}$, leaving $m$ or $m-1$:
$h(2m)=h(m)+h(m-1)$. On the other hand $u(n)=s(n+1)$ satisfies $u(2m+1)=s(2m+2)=s(m+1)=u(m)$ and
$u(2m)=s(2m+1)=s(m)+s(m+1)=u(m-1)+u(m)$, with $u(0)=1=h(0)$. By strong induction $h=u$.

(b) Expanding the product, a term $x^n$ is a choice of $e_k\in\{0,1,2\}$ in each factor with $\sum e_k2^k=n$, that is, a
hyperbinary representation of $n$. (The product converges in $\Z[[x]]$, since the $k$-th factor is $1+O(x^{2^k})$.)

(c) Let $R_k=\sum_{2^k\le n<2^{k+1}}s(n)$; $R_0=s(1)=1$. Splitting row $k+1$ into even and odd $n$,
\[
\begin{aligned}
R_{k+1}&=\sum_{2^k\le m<2^{k+1}}\bigl(s(2m)+s(2m+1)\bigr)\\
&=\sum_m\bigl(2s(m)+s(m+1)\bigr)\\
&=3R_k-s(2^k)+s(2^{k+1})=3R_k,
\end{aligned}
\]
since $s$ of a power of $2$ is $1$.

(d) $\gcd(s(0),s(1))=1$. If $\gcd(s(n),s(n+1))=1$, then $\gcd(s(2n),s(2n+1))=\gcd(s(n),s(n)+s(n+1))=1$ and
$\gcd(s(2n+1),s(2n+2))=\gcd(s(n)+s(n+1),s(n+1))=1$. By induction all consecutive pairs are coprime.

\tools Volume~1, Chapter~1 (Euclid's algorithm); Volume~4, Chapter~2 (infinite products of power series).
\leads By (d), $\frac{s(n)}{s(n+1)}$ is a fraction in lowest terms, and the sequence $\frac11,\frac12,\frac21,\frac13,\frac32,\frac23,\frac31,\dots$
lists every positive rational exactly once: it is the Calkin--Wilf enumeration of Volume~6, Exercise~1.9 (``Recounting the
rationals''). Stern introduced the sequence in 1858 in connection with the Stern--Brocot tree of Chapter~6.
\end{solution}

\begin{solution}{pb:sternfib}
\idea Hyperbinary representations are pairs of binary numbers without common digits, and such pairs are what Lucas's theorem
detects. The bound is an inequality version of the Fibonacci recursion, because of two consecutive integers one is even.

\writeup (a) Given a hyperbinary representation of $n$, let $k=\sum_{e_i=2}2^i$ and $j=\sum_{e_i=1}2^i$. Then $n=2k+j$ and $k$, $j$
have no binary digit $1$ in common. Conversely such a pair $(k,j)$ gives back the $e_i$. So $s(n+1)$ is the number of $k\ge0$
with $2k\le n$ such that $k$ and $j=n-2k$ have disjoint binary supports. By Lucas (Volume~1, Exercise~3.9), $\binom{k+j}k$ is odd if
and only if every binary digit of $k$ is at most the digit of $k+j$ in the same place, that is, if and only if there is no carry in
$k+j$, that is, if and only if $k$ and $j$ have disjoint supports. As $k+j=n-k$, this is the oddness of $\binom{n-k}k$.

(b) Let $M_k=\max\{s(n):2^k\le n\le2^{k+1}\}$, so $M_0=1=F_2$ and $M_1=\max(s(2),s(3),s(4))=2=F_3$. For $n$ in $[2^{k+1},2^{k+2}]$: if
$n=2m$ is even, $s(n)=s(m)\le M_k$. If $n=2m+1$, then $m$ and $m+1$ lie in $[2^k,2^{k+1}]$, and one of them is even, equal to $2r$
with $r\in[2^{k-1},2^k]$, so $s(n)=s(m)+s(m+1)\le M_k+M_{k-1}$. Hence $M_{k+1}\le M_k+M_{k-1}$ and $M_k\le F_{k+2}$ by induction.

(c) Let $a_k=\frac{2^{k+2}-(-1)^k}3$, so $a_1=3$, $a_2=5$, $a_3=11$, and $a_{k+1}=2a_k+\varepsilon_k$ with $\varepsilon_k=(-1)^k$ (check:
$2a_k+(-1)^k=\frac{2^{k+3}-2(-1)^k+3(-1)^k}3=a_{k+1}$). Then $s(a_{k+1})=s(a_k)+s(a_k+\varepsilon_k)$, and
$a_k+\varepsilon_k=2a_{k-1}+\varepsilon_{k-1}+\varepsilon_k=2a_{k-1}$, so $s(a_k+\varepsilon_k)=s(a_{k-1})$. Thus $s(a_{k+1})=s(a_k)+s(a_{k-1})$ with
$s(a_1)=2$, $s(a_2)=3$, and $s(a_k)=F_{k+2}$. Finally $2^k\le a_k\le2^{k+1}$.

\tools Volume~1, Exercise~3.9 (Lucas); Volume~0, Chapter~1 (Fibonacci numbers).
\leads By (a) and the identity $\sum_k\binom{n-k}k=F_{n+1}$, $s(n+1)$ counts the odd numbers on the $n$-th ``shallow diagonal''
of Pascal's triangle, whose sum is $F_{n+1}$. The second maximum point of each row, $\frac{5\cdot2^k+(-1)^k}3$ in the figure, is the
mirror image: $s(2^{k+1}+2^k-n)=s(n)$ for $2^k\le n\le2^{k+1}$, which follows from the reflection symmetry of the Stern--Brocot tree.
See Lehmer's and Northshield's papers below.
\end{solution}

\begin{solution}{pb:binarypart}
\idea Look at the smallest part: it is $1$, or every part is even and we can halve. Modulo $2$, the generating function satisfies an
equation that pins it down completely.

\writeup (a) If $n$ is odd, every partition of $n$ into powers of $2$ contains a part $1$; removing one is a bijection onto the
partitions of $n-1$ into powers of $2$, so $b(2n+1)=b(2n)$. For $2n$, the partitions with a part $1$ correspond in the same way to
those of $2n-1$, and those without have all parts even, and halving them gives the partitions of $n$. So $b(2n)=b(2n-1)+b(n)$.

(b) First proof: $b(2)=b(3)=2$. For $n\ge2$, $b(2n)=b(2n-1)+b(n)=b(2n-2)+b(n)$ is a sum of two even numbers by induction (both
$2n-2$ and $n$ are $\ge2$ and smaller than $2n$), and $b(2n+1)=b(2n)$.

Second proof: $B(x)=\sum b(n)x^n=\prod_{j\ge0}(1-x^{2^j})^{-1}$, so $(1-x)B(x)=\prod_{j\ge1}(1-x^{2^j})^{-1}=B(x^2)$. Modulo $2$, $B(x^2)=B(x)^2$, so
$(1+x)B=B^2$ in $\mathbb F_2[[x]]$, and since $B$ is invertible, $B=1+x$: all coefficients from $x^2$ on are even.

(c) Summing $b(2j)-b(2j-2)=b(j)$ for $j=1,\dots,n$ gives $b(2n)-b(0)=b(1)+\dots+b(n)$. The sequence $b$ is non-decreasing (add a part
$1$), so $b(2n)\le(n+1)b(n)$. For $n=2^k$, $b(2^k)\le\prod_{j=0}^{k-1}(2^j+1)=2^{k(k-1)/2}\prod_j(1+2^{-j})\le C\,2^{k^2/2}$. For
$2^k\le n<2^{k+1}$, monotonicity gives $\log_2b(n)\le\frac{(k+1)^2}2+O(1)\le\frac{(\log_2n)^2}2+O(\log n)$, which is the statement in
natural logarithms.

\tools Volume~4, Chapter~2 (generating functions of partitions).
\leads The upper bound has the right main term: de Bruijn (1948), following Mahler, showed that $\log b(n)\sim\frac{(\log n)^2}{2\log2}$,
with explicit lower-order terms. Churchhouse (1969) proved that $b(n)$ is never divisible by $8$ (a direct computation verifies it below $3000$), and conjectured finer congruences
for $b(2^{k+1}n)-b(2^kn)$, proved by Rødseth and Gupta.
\end{solution}

\begin{solution}{pb:fibdivides}
\idea A recurrence whose step is invertible modulo $m$ is a permutation of a finite set, so it cycles back to its start; the start
contains $F_0=0$.

\writeup (a) Let $\phi(u,v)=(v,u+v)$ on $(\Z/m\Z)^2$; it is a bijection, with inverse $(u,v)\mapsto(v-u,u)$, and
$\phi(F_n,F_{n+1})=(F_{n+1},F_{n+2})$. A bijection of a finite set has finite order, say $\phi^P=\mathrm{id}$ with $P\ge1$, so
$(F_P,F_{P+1})\equiv(F_0,F_1)=(0,1)$ and $m\mid F_P$.

(b) Let $z\ge1$ be the least index with $m\mid F_z$. Since $\gcd(F_a,F_b)=F_{\gcd(a,b)}$ (Volume~0, Exercise~1.7), if $m\mid F_n$ then
$m\mid\gcd(F_n,F_z)=F_{\gcd(n,z)}$, and $\gcd(n,z)\ge1$ (for $n\ge1$) forces $\gcd(n,z)=z$ by minimality: $z\mid n$. Conversely $F_z\mid F_{kz}$.
For $n=0$ both statements hold trivially.

(c) $J_1=J_2=1$, and $J_{n+2}\equiv J_{n+1}\pmod2$, so all $J_n$ with $n\ge1$ are odd. The step $(u,v)\mapsto(v,v+2u)$ is not injective
modulo $2$ (it sends $(0,0)$ and $(1,0)$ to the same pair), so the sequence of pairs is eventually periodic but need not return to
$(0,1)$.

\tools Volume~0, Exercise~1.7 ($\gcd$ of Fibonacci numbers); Volume~2, Chapter~1 (finite order of a permutation).
\leads The least $z(m)$ is the \emph{rank of apparition}; it is known that $z(m)\le2m$, with equality for $m\ge2$ exactly when
$m=6\cdot5^k$ (a direct computation verifies this below $2000$). For primes, Problem~\ref{pb:fibmodp} shows $z(p)\mid p-\bigl(\frac5p\bigr)$. The same argument shows that any integer recurrence
whose last coefficient is $\pm1$ and which starts with $0$ hits every modulus.
\end{solution}

\begin{solution}{pb:thuemorse}
\idea The rules $t_{2n}=t_n$ and $t_{2n+1}=1+t_n$ say that $T(x)$ is expressed through $T(x^2)$, which modulo $2$ is $T(x)^2$. Periodicity
is incompatible with these rules, because a period could be halved until it becomes odd, and an odd period contradicts the
alternation $t_{2n+1}\ne t_{2n}$.

\writeup (a) Since $s(2n)=s(n)$ and $s(2n+1)=s(n)+1$,
\[
T(x)=\sum_nt_nx^{2n}+\sum_n(1+t_n)x^{2n+1}=(1+x)T(x^2)+\frac x{1-x^2}.
\]
In $\mathbb F_2[[x]]$, $T(x^2)=T(x)^2$ and $1-x^2=(1+x)^2$. So $T=(1+x)T^2+\frac x{(1+x)^2}$; multiply by $(1+x)^2$.

(b) Suppose $t_{n+P}=t_n$ for all $n\ge N$, with $P\ge1$ minimal. If $P$ is even, then for $n\ge N$, $t_{n+P/2}=t_{2n+P}=t_{2n}=t_n$, so $P/2$ is a
period from $N$ on, contradicting minimality. So $P$ is odd. For $n\ge N$,
\[
t_{2n+P}=t_{2n}=t_n\quad\text{and}\quad t_{2n+P}=t_{2(n+\frac{P-1}2)+1}=1+t_{n+\frac{P-1}2},
\]
so $t_{n+\frac{P-1}2}=1+t_n$. Similarly $t_{2n+1+P}=t_{2n+1}=1+t_n$ and $t_{2n+1+P}=t_{2(n+\frac{P+1}2)}=t_{n+\frac{P+1}2}$, so $t_{n+\frac{P+1}2}=1+t_n$.
Hence $t_{m+1}=t_m$ for all large $m$: the sequence is eventually constant. But $t_{2n+1}=1+t_{2n}$ for every $n$. Contradiction.

(c) By (a), $T$ is a root of $(1+x)^3Y^2+(1+x)^2Y+x$, so it is algebraic of degree at most $2$. If $T$ were rational, its coefficients
would satisfy a linear recurrence from some point on, and be eventually periodic as in Problem~\ref{pb:kolberg}; this contradicts
(b). So the degree is exactly $2$.

\tools Volume~4, Theorem~2.5 (rational series and recurrences); Volume~2, Chapter~9 (the Frobenius map in characteristic $p$).
\leads \emph{Christol's theorem} (1979): a power series over $\mathbb F_p$ is algebraic over $\mathbb F_p(x)$ if and only if its sequence
of coefficients is \emph{$p$-automatic}, that is, $a_n$ can be computed by a finite automaton reading the base-$p$ digits of $n$. Thue--Morse,
the Catalan numbers modulo $2$ and the central binomial coefficients modulo $p$ (Problem~\ref{pb:centralp}) are examples. See
Allouche and Shallit's book.
\end{solution}

\begin{solution}{pb:orderedfact}
\idea An ordered factorisation into $k$ factors $>1$ is counted by the Dirichlet series $(\zeta(s)-1)^k$, and the geometric series in
$\zeta-1$ converges exactly when $\zeta(s)<2$. The number $\rho$ where $\zeta(\rho)=2$ is the natural growth exponent, and a direct
induction shows $H(n)\le n^\rho$.

\writeup (a) For real $s>1$, $\zeta(s)-1=\sum_{n\ge2}n^{-s}$, so $(\zeta(s)-1)^k=\sum_nH_k(n)n^{-s}$, where $H_k(n)$ counts the ordered
factorisations of $n$ into exactly $k$ factors $>1$. Since $H=\sum_kH_k$ (with $H_0=\mathbf1_{n=1}$) and all terms are non-negative,
$\sum_nH(n)n^{-s}=\sum_k(\zeta(s)-1)^k$, which converges to $\frac1{2-\zeta(s)}$ when $\zeta(s)<2$, that is, when $s>\rho$ ($\zeta$ is decreasing).

(b) An ordered factorisation of $p^a$ is a composition $a=a_1+\dots+a_k$ with $a_i\ge1$, and there are $2^{a-1}$ compositions of $a$.
For $n=p_1\cdots p_k$, an ordered factorisation is an ordered partition of $\{p_1,\dots,p_k\}$ into non-empty blocks: $1,3,13,75$
for $k=1,2,3,4$ (the Fubini numbers; for $k=3$: $1$ factorisation with one factor, $6$ with two, $6$ with three).

(c) The first factor of an ordered factorisation of $n\ge2$ is a divisor $e>1$, and the rest is an ordered factorisation of $n/e$, so
$H(n)=\sum_{d\mid n,\,d<n}H(d)$. Assume $H(d)\le d^\rho$ for $d<n$ ($H(1)=1$). Then
\[
H(n)\le\sum_{d\mid n,\,d<n}d^\rho=n^\rho\sum_{e\mid n,\,e>1}e^{-\rho}<n^\rho\sum_{e\ge2}e^{-\rho}=n^\rho(\zeta(\rho)-1)=n^\rho.
\]

(d) For $\sigma>\rho$ and $n\le x$, $1\le(x/n)^\sigma$, so $\sum_{n\le x}H(n)\le x^\sigma\sum_nH(n)n^{-\sigma}=\frac{x^\sigma}{2-\zeta(\sigma)}$.

\tools Volume~5, Chapter~2 (Dirichlet series, convolution); Volume~4, Chapter~10 (the function $\zeta$ on $]1,\infty[$).
\leads Kalmár proved in 1931 that $\sum_{n\le x}H(n)\sim-\frac{x^\rho}{\rho\zeta'(\rho)}=0.3181\ldots x^\rho$, by a Tauberian theorem
applied at the simple pole of $\frac1{2-\zeta(s)}$ at $s=\rho$; a direct computation verifies it. The bound (c) is due to Chor, Lemke and
Mador (2000), and the exponent $\rho$ cannot be lowered: Hille showed in 1936 that $H(n)\ge n^{\rho-\varepsilon}$ for infinitely many $n$.
\end{solution}

\begin{solution}{pb:dsquared}
\idea Euler products turn the identity into one about power series in a single prime, and the asymptotics follow from the known
average of $\tau_4$ by a convolution with a sparse function.

\writeup (a) Both sides are Euler products of multiplicative functions, absolutely convergent for $s>1$. With $x=p^{-s}$, the
factor at $p$ on the left is $\sum_k(k+1)^2x^k=\frac{1+x}{(1-x)^3}$ (differentiate $\sum x^{k+1}$ twice), and on the right it is
$\frac{1-x^2}{(1-x)^4}=\frac{1+x}{(1-x)^3}$.

(b) By (a), $d^2=\tau_4*g$, where $\tau_4=\mathbf1^{*4}$ has Dirichlet series $\zeta^4$ and $g(m^2)=\mu(m)$, $g=0$ off the squares (its
series is $\frac1{\zeta(2s)}$). So $\sum_{n\le x}d(n)^2=\sum_{m\le\sqrt x}\mu(m)T(x/m^2)$ with $T(y)=\sum_{n\le y}\tau_4(n)$. By Volume~5,
Exercise~2.7, $T(y)=\frac16y\log^3y+O(y\log^2y)$ for $y\ge3$, and trivially $T(y)=O(y)$ for $1\le y<3$. Hence
\[
\sum_{n\le x}d(n)^2=\frac x6\sum_{m\le\sqrt x}\frac{\mu(m)}{m^2}\Bigl(\log\frac x{m^2}\Bigr)^3+O\Bigl(\sum_{m\le\sqrt x}\frac x{m^2}\log^2x\Bigr).
\]
The error is $O(x\log^2x)$. Expanding $(\log x-2\log m)^3=\log^3x+O(\log^2x\,(1+\log m)^3)$ and using $\sum_m\frac{(1+\log m)^3}{m^2}<\infty$,
the main term is $\frac x6\log^3x\sum_{m\le\sqrt x}\frac{\mu(m)}{m^2}+O(x\log^2x)$. Finally
$\sum_{m\le\sqrt x}\frac{\mu(m)}{m^2}=\frac1{\zeta(2)}+O(x^{-1/2})=\frac6{\pi^2}+O(x^{-1/2})$, and $\frac16\cdot\frac6{\pi^2}=\frac1{\pi^2}$.

\tools Volume~5, Exercise~2.7 (average of $\tau_k$); Volume~1, Chapter~4 (Möbius, multiplicative functions); Volume~4, Chapter~10
($\zeta(2)=\frac{\pi^2}6$).
\leads Ramanujan stated (a) in 1916, together with similar identities such as $\sum\sigma_a(n)\sigma_b(n)n^{-s}=\frac{\zeta(s)\zeta(s-a)\zeta(s-b)\zeta(s-a-b)}{\zeta(2s-a-b)}$.
The true asymptotic is $xP(\log x)$ with $P$ a cubic, and the convergence of the leading term is slow: one finds numerically
$\frac{\sum d(n)^2}{x\log^3x}\approx0.17$ at $x=10^5$, against the limit $0.101$.
\end{solution}

\begin{solution}{pb:powerful}
\idea A powerful number has every exponent $\ge2$, and every such exponent is $2u+3v$ with $v\in\{0,1\}$. So powerful numbers are
squares times cubes of squarefree numbers, and counting them is counting lattice points under a curve.

\writeup (a) Let $n=\prod p^{e_p}$ with all $e_p\ge2$. Put $v_p=1$ if $e_p$ is odd and $v_p=0$ otherwise, and $u_p=\frac{e_p-3v_p}2\ge0$. Then
$n=a^2b^3$ with $a=\prod p^{u_p}$ and $b=\prod p^{v_p}$ squarefree. Conversely $n=a^2b^3$ with $b$ squarefree forces $v_p\equiv e_p\pmod2$,
so the decomposition is unique, and every $a^2b^3$ is powerful.

(b) By (a) the count is
\[
\sum_{b\le x^{1/3}}\mu^2(b)\Bigl\lfloor\sqrt{x/b^3}\Bigr\rfloor=\sqrt x\sum_{b\le x^{1/3}}\frac{\mu^2(b)}{b^{3/2}}+O(x^{1/3}).
\]
Since $\sum_{b>y}b^{-3/2}=O(y^{-1/2})$, completing the sum costs $O(\sqrt x\,x^{-1/6})=O(x^{1/3})$, and
$\sum_b\mu^2(b)b^{-3/2}=\frac{\zeta(3/2)}{\zeta(3)}$ because $\sum\mu^2(b)b^{-s}=\prod_p(1+p^{-s})=\frac{\zeta(s)}{\zeta(2s)}$. For the Dirichlet series,
the factor at $p$ is, with $y=p^{-s}$,
\[
1+\sum_{e\ge2}y^e=\frac{1-y+y^2}{1-y},\qquad\text{and}\qquad\frac{1-y^6}{(1-y^2)(1-y^3)}=\frac{1+y^3}{1-y^2}=\frac{(1+y)(1-y+y^2)}{(1-y)(1+y)},
\]
the same.

(c) Pell's equation $X^2-8Y^2=1$ has infinitely many solutions (Volume~1, Theorem~8.7), for instance $(3,1)$ and $(17,6)$, the
others coming from powers of $3+\sqrt8$. For each, $8Y^2=2^3Y^2$ is powerful (every prime of $Y$ appears at least squared, and $2$ at
least cubed) and $8Y^2+1=X^2$ is a square. This gives $(8,9)$ and $(288,289)$.

\tools Volume~1, Theorem~8.7 (Pell's equation); Volume~5, Chapter~2 (counting with Dirichlet series).
\leads Erdős conjectured that there are no three consecutive powerful numbers; this is still open (and would follow from the
$abc$ conjecture). The next term of the asymptotic expansion is $\frac{\zeta(2/3)}{\zeta(2)}x^{1/3}$, negative since $\zeta(2/3)<0$, which is why
the numerical counts lie below $2.173\sqrt x$.
\end{solution}

\begin{solution}{pb:abelian}
\idea The structure theorem reduces abelian groups of order $p^e$ to partitions of $e$, and the generating function of partitions
becomes an Euler product of zeta values. The average is governed by the simple pole of $\zeta(s)$ at $1$; the other factors are a
convergent correction.

\writeup (a) A finite abelian group is the direct product of its Sylow subgroups, and an abelian $p$-group of order $p^e$ is
$\prod_i\Z/p^{\lambda_i}\Z$ for a unique partition $\lambda$ of $e$ (Volume~2, Corollary~7.7). So $A$ is multiplicative with $A(p^e)=p(e)$. Its
Euler factor is $\sum_ep(e)p^{-es}=\prod_{k\ge1}(1-p^{-ks})^{-1}$, and the product over $p$ is $\prod_k\zeta(ks)$ (all products converge
absolutely for $s>1$).

(b) Write $\prod_k\zeta(ks)=\zeta(s)G(s)$ with $G(s)=\prod_{k\ge2}\zeta(ks)=\sum_mg(m)m^{-s}$. Here $g\ge0$, and $g$ is supported on the
powerful numbers (its Euler factor at $p$ has no term $p^{-s}$). Then $A=\mathbf1*g$ and
\[
\sum_{n\le x}A(n)=\sum_{m\le x}g(m)\Bigl\lfloor\frac xm\Bigr\rfloor=x\sum_mg(m)m^{-1}-x\sum_{m>x}\frac{g(m)}m+O\Bigl(\sum_{m\le x}g(m)\Bigr).
\]
Write $G(s)=\zeta(2s)K(s)$ with $K(s)=\prod_{k\ge3}\zeta(ks)=\sum_jh(j)j^{-s}$, $h\ge0$, convergent at every $s>\frac13$. Then
$\sum_{m\le y}g(m)=\sum_jh(j)\lfloor\sqrt{y/j}\rfloor\le\sqrt y\sum_jh(j)j^{-1/2}=K(\tfrac12)\sqrt y$. Writing $G(y)=\sum_{m\le y}g(m)$, partial summation gives
$\sum_{m>x}\frac{g(m)}m=-\frac{G(x)}x+\int_x^\infty\frac{G(y)}{y^2}dy\le\int_x^\infty\frac{K(\frac12)}{y^{3/2}}dy=\frac{2K(\frac12)}{\sqrt x}$. So the sum is $G(1)x+O(\sqrt x)$, with
$G(1)=\prod_{k\ge2}\zeta(k)=2.2948\ldots$

\tools Volume~2, Corollary~7.7 (finitely generated abelian groups); Volume~4, Proposition~2.11 (generating function of $p(n)$);
Volume~5, Chapter~2 (convolutions and partial summation).
\leads Erdős and Szekeres proved (b) in 1934. The truth is $\sum_{n\le x}A(n)=c_1x+c_2x^{1/2}+c_3x^{1/3}+o(x^{1/3})$ with $c_k=\prod_{j\ne k}\zeta(\frac jk)$, so
$c_2<0$ because $\zeta(\frac12)<0$; the best exponent in the error term is an open problem. The individual values of $A(n)$ can be large: $A(2^{e})=p(e)$ grows like
$e^{c\sqrt e}$.
\end{solution}

\section*{Map of the chapter}

\begin{center}
\small
\adjustbox{max width=\linewidth}{\begin{tabular}{lll}
Problem & Topic & Volumes used\\\hline
\ref{pb:parts123} & partitions into $1,2,3$; integer triangles & 4\\
\ref{pb:selfconj} & parity of $p(n)$, self-conjugate partitions & 4\\
\ref{pb:glaisherthm} & Glaisher's theorem and bijection & 4\\
\ref{pb:stanley} & Stanley's and Elder's theorems & 0\\
\ref{pb:eulerlog} & logarithmic derivative, $np(n)=\sum\sigma(k)p(n-k)$ & 2, 4\\
\ref{pb:kolberg} & both parities of $p(n)$ & 4\\
\ref{pb:ramanujan5} & triple product, $p(5n+4)\equiv0\pmod5$ & 1, 4\\
\ref{pb:kummer} & Kummer's theorem, lcm bound & 1\\
\ref{pb:catalanodd} & odd Catalan numbers & 3, 4\\
\ref{pb:pascalthree} & Pascal's triangle modulo $3$ & 1\\
\ref{pb:pascalfour} & Pascal's triangle modulo $4$ & 1\\
\ref{pb:centralp} & $\binom{2n}n$ modulo $p$ & 1, 3\\
\ref{pb:touchard} & Touchard's congruence for Bell numbers & 2, 4\\
\ref{pb:stern} & Stern's sequence, hyperbinary expansions & 1, 4\\
\ref{pb:sternfib} & Stern, Pascal and Fibonacci & 0, 1\\
\ref{pb:binarypart} & binary partitions & 4\\
\ref{pb:fibdivides} & every integer divides a Fibonacci number & 0, 2\\
\ref{pb:thuemorse} & Thue--Morse is algebraic over $\mathbb F_2$ & 2, 4\\
\ref{pb:orderedfact} & ordered factorisations & 4, 5\\
\ref{pb:dsquared} & $\sum d(n)^2$ & 1, 4, 5\\
\ref{pb:powerful} & powerful numbers & 1, 5\\
\ref{pb:abelian} & counting abelian groups & 2, 4, 5\\
\end{tabular}}
\end{center}

\section*{Sources and further reading}

\begin{itemize}
\item G.~H.~Hardy, Srinivasa Ramanujan, \emph{Proceedings of the London Mathematical Society} (2) 19 (1921), xl--lviii (the epigraph:
``It was Mr Littlewood (I believe) who remarked that `every positive integer was one of his personal friends'\,''); G.~H.~Hardy,
\emph{Ramanujan: Twelve Lectures on Subjects Suggested by His Life and Work}, Cambridge, 1940, Lecture~VIII.
\item G.~H.~Hardy and S.~Ramanujan, Asymptotic formulae in combinatory analysis, \emph{Proceedings of the London Mathematical
Society} (2) 17 (1918), 75--115 (MacMahon's table of $p(n)$, $n\le200$, and the check of $p(200)$).
\item S.~Ramanujan, Some properties of $p(n)$, the number of partitions of $n$, \emph{Proceedings of the Cambridge Philosophical
Society} 19 (1919), 207--210; S.~Ramanujan, Some formulae in the analytic theory of numbers, \emph{Messenger of Mathematics} 45
(1916), 81--84 (Problem~\ref{pb:dsquared}).
\item F.~J.~Dyson, Some guesses in the theory of partitions, \emph{Eureka} 8 (1944), 10--15; A.~O.~L.~Atkin and H.~P.~F.~Swinnerton-Dyer,
Some properties of partitions, \emph{Proceedings of the London Mathematical Society} (3) 4 (1954), 84--106; G.~E.~Andrews and
F.~G.~Garvan, Dyson's crank of a partition, \emph{Bulletin of the AMS} 18 (1988), 167--171.
\item G.~E.~Andrews, \emph{The Theory of Partitions}, Addison-Wesley, 1976 (Cambridge reprint 1998): the triple product, Glaisher's
theorem, self-conjugate partitions.
\item K.~Ono, Distribution of the partition function modulo $m$, \emph{Annals of Mathematics} 151 (2000), 293--307.
\item O.~Kolberg, Note on the parity of the partition function, \emph{Mathematica Scandinavica} 7 (1959), 377--378.
\item K.~Ono, Parity of the partition function in arithmetic progressions, \emph{Journal für die reine und angewandte Mathematik}
472 (1996), 1--15; C.-S.~Radu, A proof of Subbarao's conjecture, \emph{Journal für die reine und angewandte Mathematik} 672 (2012),
161--175; J.-L.~Nicolas, I.~Z.~Ruzsa and A.~Sárközy, On the parity of additive representation functions, \emph{Journal of Number
Theory} 73 (1998), 292--317.
\item J.~W.~L.~Glaisher, A theorem in partitions, \emph{Messenger of Mathematics} 12 (1883), 158--170.
\item R.~Honsberger, \emph{Mathematical Gems~III}, MAA, 1985 (Elder's theorem).
\item E.~E.~Kummer, Über die Ergänzungssätze zu den allgemeinen Reciprocitätsgesetzen, \emph{Journal für die reine und angewandte
Mathematik} 44 (1852), 93--146.
\item A.~Granville, Zaphod Beeblebrox's brain and the fifty-ninth row of Pascal's triangle, \emph{American Mathematical Monthly} 99
(1992), 318--331. R.~Garfield and H.~S.~Wilf, The distribution of the binomial coefficients modulo $p$, \emph{Journal of Number Theory}
41 (1992), 1--5; I.~M.~Gessel, Some congruences for Apéry numbers, \emph{Journal of Number Theory} 14 (1982), 362--368.
\item J.~Touchard, Propriétés arithmétiques de certains nombres récurrents, \emph{Annales de la Société Scientifique de Bruxelles}
A 53 (1933), 21--31.
\item M.~A.~Stern, Über eine zahlentheoretische Funktion, \emph{Journal für die reine und angewandte Mathematik} 55 (1858),
193--220; D.~H.~Lehmer, On Stern's diatomic series, \emph{American Mathematical Monthly} 36 (1929), 59--67; S.~Northshield, Stern's
diatomic sequence $0,1,1,2,1,3,2,3,1,4,\dots$, \emph{American Mathematical Monthly} 117 (2010), 581--598.
\item N.~G.~de Bruijn, On Mahler's partition problem, \emph{Indagationes Mathematicae} 10 (1948), 210--220; R.~F.~Churchhouse,
Congruence properties of the binary partition function, \emph{Proceedings of the Cambridge Philosophical Society} 66 (1969),
371--376.
\item J.-P.~Allouche and J.~Shallit, \emph{Automatic Sequences}, Cambridge, 2003; G.~Christol, Ensembles presque périodiques
$k$-reconnaissables, \emph{Theoretical Computer Science} 9 (1979), 141--145.
\item L.~Kalmár, Über die mittlere Anzahl der Produktdarstellungen der Zahlen, \emph{Acta Scientiarum Mathematicarum (Szeged)} 5
(1931), 95--107; B.~Chor, P.~Lemke and Z.~Mador, On the number of ordered factorizations of natural numbers, \emph{Discrete
Mathematics} 214 (2000), 123--133. E.~Hille, A problem in ``Factorisatio Numerorum'', \emph{Acta Arithmetica} 2 (1936), 134--144.
\item P.~Erdős and G.~Szekeres, Über die Anzahl der Abelschen Gruppen gegebener Ordnung und über ein verwandtes
zahlentheoretisches Problem, \emph{Acta Scientiarum Mathematicarum (Szeged)} 7 (1934), 95--102; S.~W.~Golomb, Powerful numbers,
\emph{American Mathematical Monthly} 77 (1970), 848--852.
\item The On-Line Encyclopedia of Integer Sequences, \texttt{oeis.org}, sequences A000041 ($p(n)$), A002487 (Stern), A005044
(Alcuin), A018819 (binary partitions), A074206 (ordered factorisations).
\end{itemize}
